The core ontology architecture bridges the gap between cyber-physical energy systems and deep learning signal processing domain knowledge. It categorizes the application setting into five main conceptual pillars rooted under the ontological hierarchy:

\begin{itemize}
    \item \textbf{Physical Subsystem (\texttt{pv:PhysicalComponent}):} Represents the hardware infrastructure of the micro-grid, including the solar array, electrochemical storage, consumer loads, and the macro grid tie.
    \item \textbf{Cognitive Layer (\texttt{pv:ControlAgent}):} The artificial intelligence system responsible for autonomous grid optimization, load balancing, dispatch scheduling, and proactive vulnerability mitigation.
    \item \textbf{Signal Processing Layer (\texttt{pv:Signal}, \texttt{pv:PatchedSignal}):} Tracks the raw telemetry inputs (both time-domain features and frequency-domain components) along with their tokenized representation configurations (patch dimensions, stride values) tailored for the Chronos foundation architecture.
    \item \textbf{Environmental Tracking (\texttt{pv:WeatherObservation}):} Captures concurrent meteorological snapshots (air temperature, wind vectors) influencing structural physics and electrical efficiency.
    \item \textbf{Inferred Operational States:} Dynamic metadata classified into generation levels (\texttt{pv:GenerationState}), systemic disruptions (\texttt{pv:AnomalyState}), control directives (\texttt{pv:AgentAction}), and signal degradation hazards (\texttt{pv:AliasingEvent}).
\end{itemize}

\section{Core Concept Taxonomy}
The formal classes defined within the schema are explicitly detailed below:
\begin{description}
    \item[\texttt{pv:PhotovoltaicArray}] The primary renewable energy asset. Characterized by panel count and nameplate nominal capacity. Source of telemetry streams; associated with weather conditions and subject to environmental/electrical anomalies.
    \item[\texttt{pv:BatteryStorage}] Local electrochemical reservoir capable of dynamic grid connection/isolation. Characterized by capacity and instantaneous State of Charge (SoC).
    \item[\texttt{pv:InternalLoad}] Represents internal industrial or residential consumption demands. Subject to adjustable, demand-side management interventions.
    \item[\texttt{pv:NationalGrid}] The external electrical macro-grid interface, acting as a bidirectional infinite bus for excess export or deficit imports.
    \item[\texttt{pv:Signal}] Time-series measurement stream containing global horizontal irradiance ($G_h$), tilted surface irradiance ($G_{\text{tilt}}$), instantaneous power output ($P_{\text{pv}}$), and analytical frequency descriptors (sampling rate, fundamental frequency, amplitude, phase).
    \item[\texttt{pv:PatchedSignal}] The Chronos-specific tokenized structure derived from a raw signal. Governed by a selected window chunking mechanism ($patchSize$, $stride$), establishing a structural $overlapRatio$ and a downstream $patchInducedFrequency$.
\end{description}

\section{Semantic Relations \& Properties}

\subsection{Object Properties (Entity-to-Entity Relationships)}
The physical, functional, and inferred linkages across concepts are formalized using the following object properties:
\begin{itemize}
    \item $\langle\texttt{pv:PhotovoltaicArray}\rangle \ \texttt{pv:measures} \ \langle\texttt{pv:Signal}\rangle$
    \item $\langle\texttt{pv:PhotovoltaicArray}\rangle \ \texttt{pv:hasWeatherObservation} \ \langle\texttt{pv:WeatherObservation}\rangle$
    \item $\langle\texttt{pv:PhotovoltaicArray}\rangle \ \texttt{pv:hasGenerationState} \ \langle\texttt{pv:GenerationState}\rangle$ \quad \textit{(Inferred via SWRL)}
    \item $\langle\texttt{pv:PhotovoltaicArray}\rangle \ \texttt{pv:hasAnomaly} \ \langle\texttt{pv:AnomalyState}\rangle$ \quad \textit{(Inferred via SWRL)}
    \item $\langle\texttt{pv:PhotovoltaicArray}\rangle \ \texttt{pv:triggersAction} \ \langle\texttt{pv:AgentAction}\rangle$ \quad \textit{(Inferred via SWRL)}
    \item $\langle\texttt{pv:Signal}\rangle \ \texttt{pv:tokenizes} \ \langle\texttt{pv:PatchedSignal}\rangle$
    \item $\langle\texttt{pv:PatchedSignal}\rangle \ \texttt{pv:exhibitsAliasing} \ \langle\texttt{pv:AliasingEvent}\rangle$ \quad \textit{(Inferred via SWRL)}
    \item $\langle\texttt{pv:BatteryStorage}\rangle \ \texttt{pv:buffers} \ \langle\texttt{pv:PhotovoltaicArray}\rangle$
    \item $\langle\texttt{pv:PhysicalComponent}\rangle \ \texttt{pv:connectedTo} \ \langle\texttt{pv:PhysicalComponent}\rangle$ \quad \textit{(Symmetric Property)}
\end{itemize}

\subsection{Data Properties (Entity Attributes)}
Key quantitative attributes defining system conditions include:
\begin{itemize}
    \item \textbf{Electrical/Solar:} \texttt{pv:pvPowerW} ($P_{\text{pv}}$), \texttt{pv:nominalPowerW} ($P_{\text{nom}}$), \texttt{pv:tiltedIrradianceWm2} ($G_{\text{tilt}}$), \texttt{pv:deltaPowerW} ($\Delta P$).
    \item \textbf{Meteorological:} \texttt{pv:airTemperatureC} ($T_{\text{air}}$), \texttt{pv:windSpeedMs} ($W_s$).
    \item \textbf{Tokenization Architecture:} \texttt{pv:patchSize} ($N_p$), \texttt{pv:stride} ($S$), \texttt{pv:overlapRatio} ($\gamma$), \texttt{pv:patchInducedFrequencyHz} ($f_p$).
    \item \textbf{Signal Descriptors:} \texttt{pv:samplingFrequencyHz} ($f_s$), \texttt{pv:fundamentalFrequencyHz} ($f_0$), \texttt{pv:phaseRad} ($\phi$).
\end{itemize}

\section{Knowledge-Driven Inference System (SWRL Rules)}
Real-time diagnostic, classification, and mitigation logic are automated using a rule-based engine. The standard logical form of these production rules is detailed below:

\subsection{Environmental Context Extraction}
\begin{equation}
\begin{aligned}
\texttt{PhotovoltaicArray}(?pv) &\wedge \texttt{measures}(?pv, ?s) \wedge \texttt{tiltedIrradianceWm2}(?s, ?g) \wedge (?g \ge 800.0) \\
&\longrightarrow \texttt{hasGenerationState}(?pv, \texttt{ClearSky})
\end{aligned}
\end{equation}

\begin{equation}
\begin{aligned}
\texttt{PhotovoltaicArray}(?pv) &\wedge \texttt{measures}(?pv, ?s) \wedge \texttt{tiltedIrradianceWm2}(?s, ?g) \\
&\wedge (?g \ge 150.0) \wedge (?g < 800.0) \longrightarrow \texttt{hasGenerationState}(?pv, \texttt{PartialCloud})
\end{aligned}
\end{equation}

\begin{equation}
\begin{aligned}
\texttt{PhotovoltaicArray}(?pv) &\wedge \texttt{measures}(?pv, ?s) \wedge \texttt{tiltedIrradianceWm2}(?s, ?g) \\
&\wedge (?g > 0.0) \wedge (?g < 150.0) \longrightarrow \texttt{hasGenerationState}(?pv, \texttt{Overcast})
\end{aligned}
\end{equation}

\begin{equation}
\begin{aligned}
\texttt{PhotovoltaicArray}(?pv) &\wedge \texttt{measures}(?pv, ?s) \wedge \texttt{tiltedIrradianceWm2}(?s, ?g) \wedge (?g \le 0.0) \\
&\longrightarrow \texttt{hasGenerationState}(?pv, \texttt{Nighttime})
\end{aligned}
\end{equation}

\subsection{Cyber-Physical Anomaly Diagnostics}
\begin{itemize}
    \item \textbf{Inverter Clipping Anomaly (R5):}
    \begin{equation}
    \begin{aligned}
    \texttt{PhotovoltaicArray}(?pv) &\wedge \texttt{measures}(?pv, ?s) \wedge \texttt{pvPowerW}(?s, ?p) \wedge \texttt{nominalPowerW}(?pv, ?nom) \\
    &\wedge (?p \ge ?nom) \wedge \texttt{tiltedIrradianceWm2}(?s, ?g) \wedge (?g > 800.0) \\
    &\longrightarrow \texttt{hasAnomaly}(?pv, \texttt{Clipping})
    \end{aligned}
    \end{equation}
    
    \item \textbf{Inverter Trip Anomaly (R6):}
    \begin{equation}
    \begin{aligned}
    \texttt{PhotovoltaicArray}(?pv) &\wedge \texttt{measures}(?pv, ?s) \wedge \texttt{pvPowerW}(?s, ?p) \wedge (?p \le 0.0) \\
    &\wedge \texttt{tiltedIrradianceWm2}(?s, ?g) \wedge (?g > 200.0) \longrightarrow \texttt{hasAnomaly}(?pv, \texttt{InverterTrip})
    \end{aligned}
    \end{equation}

    \item \textbf{Ramp Event Anomaly (R7):}
    \begin{equation}
    \texttt{PhotovoltaicArray}(?pv) \wedge \texttt{measures}(?pv, ?s) \wedge \texttt{deltaPowerW}(?s, ?dp) \wedge (?dp > 500.0) \longrightarrow \texttt{hasAnomaly}(?pv, \texttt{Ramp})
    \end{equation}

    \item \textbf{Soiling Performance Deficit Anomaly (R8):}
    \begin{equation}
    \begin{aligned}
    \texttt{PhotovoltaicArray}(?pv) &\wedge \texttt{measures}(?pv, ?s) \wedge \texttt{tiltedIrradianceWm2}(?s, ?g) \wedge (?g > 400.0) \\
    &\wedge \texttt{pvPowerW}(?s, ?p) \wedge (?p > 100.0) \wedge \texttt{swrlb:divide}(?r, ?p, ?g) \wedge (?r < 0.3) \\
    &\longrightarrow \texttt{hasAnomaly}(?pv, \texttt{Soiling})
    \end{aligned}
    \end{equation}

    \item \textbf{Thermal Overheating Derating Anomaly (R9):}
    \begin{equation}
    \begin{aligned}
    \texttt{PhotovoltaicArray}(?pv) &\wedge \texttt{hasWeatherObservation}(?pv, ?w) \wedge \texttt{airTemperatureC}(?w, ?t) \wedge (?t > 35.0) \\
    &\wedge \texttt{measures}(?pv, ?s) \wedge \texttt{tiltedIrradianceWm2}(?s, ?g) \wedge (?g > 600.0) \\
    &\wedge \texttt{pvPowerW}(?s, ?p) \wedge \texttt{swrlb:divide}(?r, ?p, ?g) \wedge (?r < 0.4) \\
    &\longrightarrow \texttt{hasAnomaly}(?pv, \texttt{Overheating})
    \end{aligned}
    \end{equation}
\end{itemize}

\subsection{Autonomous Closed-Loop Control Actuation}
Based on inferred operating profiles, actions are dynamically mapped to handle grid stresses:
\begin{align}
\texttt{pv:hasAnomaly}(?pv, \texttt{pv:Ramp}) &\longrightarrow \texttt{pv:triggersAction}(?pv, \texttt{pv:RampAlert}) \\
\texttt{pv:hasAnomaly}(?pv, \texttt{pv:Clipping}) &\vee \texttt{pv:hasAnomaly}(?pv, \texttt{pv:Overheating}) \\ \longrightarrow \texttt{pv:triggersAction}(?pv, \texttt{pv:CurtailGeneration}) \\
\texttt{pv:hasAnomaly}(?pv, \texttt{pv:InverterTrip}) &\vee \texttt{pv:hasGenerationState}(?pv, \texttt{pv:PartialCloud}) \\ \longrightarrow \texttt{pv:triggersAction}(?pv, \texttt{pv:BatteryDispatch}) \\
\texttt{pv:hasGenerationState}(?pv, \texttt{pv:Overcast}) &\vee \texttt{pv:hasGenerationState}(?pv, \texttt{pv:Nighttime}) \\ \longrightarrow \texttt{pv:triggersAction}(?pv, \texttt{pv:LoadAdjustment})
\end{align}

\section{Chronos Structural Aliasing Hypotheses}
The key scientific contribution of this architecture lies in identifying information loss bounds during time-series patching. The model detects structural blind spots via frequency domain reasoning:

\subsection{Hypothesis H1 \& H1b: Patch-Induced Frequency Blind Spots}
The patch tokenization process introduces an artificial operating frequency profile $f_p$, mathematically defined by the relationship:
\begin{equation}
f_p = \frac{f_s}{S}
\end{equation}
Where $f_s$ represents the raw Nyquist sampling rate and $S$ denotes the token step stride. 

\begin{itemize}
    \item \textbf{Primary Blind Spot (H1):} If the continuous signal's structural fundamental frequency ($f_0$, such as the daily diurnal solar cycle) aligns directly with the token boundaries, a tracking blind spot forms at the fundamental harmonic region ($n=1$):
    \begin{equation}
    \begin{aligned}
    \texttt{Signal}(?s) &\wedge \texttt{tokenizes}(?s, ?ps) \wedge \texttt{fundamentalFrequencyHz}(?s, ?f_0) \wedge \texttt{patchInducedFrequencyHz}(?ps, ?f_p) \\
    &\wedge \texttt{swrlb:divide}(?r, ?f_0, ?f_p) \wedge (?r \ge 0.9) \wedge (?r \le 1.1) \\
    &\longrightarrow \texttt{exhibitsAliasing}(?ps, \texttt{H1\_BlindSpot})
    \end{aligned}
    \end{equation}
    
    \item \textbf{Secondary Harmonic Blind Spot (H1b):} Re-occurs at the second harmonic ($n=2$) boundary layer:
    \begin{equation}
    \begin{aligned}
    \texttt{Signal}(?s) &\wedge \texttt{tokenizes}(?s, ?ps) \wedge \texttt{fundamentalFrequencyHz}(?s, ?f_0) \wedge \texttt{patchInducedFrequencyHz}(?ps, ?f_p) \\
    &\wedge \texttt{swrlb:divide}(?r, ?f_0, ?f_p) \wedge (?r \ge 1.85) \wedge (?r \le 2.15) \\
    &\longrightarrow \texttt{exhibitsAliasing}(?ps, \texttt{H1b\_BlindSpot})
    \end{aligned}
    \end{equation}
\end{itemize}

\subsection{Hypothesis H2: Phase Aliasing Effect}
Vulnerabilities appear when temporal offsets misalign with chunk window anchors, corrupting early representation extraction features ($0 < \phi < S$):
\begin{equation}
\begin{aligned}
\texttt{Signal}(?s) &\wedge \texttt{tokenizes}(?s, ?ps) \wedge \texttt{phaseRad}(?s, ?\phi) \wedge \texttt{stride}(?ps, ?S) \\
&\wedge (?\phi > 0.0) \wedge (?\phi < ?S) \longrightarrow \texttt{exhibitsAliasing}(?ps, \texttt{H2\_Phase})
\end{aligned}
\end{equation}

\subsection{Hypothesis H3: Overlap Information Distortion}
When the overlap configuration ratio ($\gamma = \frac{N_p - S}{N_p}$) drops below a critical information-theoretic bound ($\gamma < 0.5$), consecutive contextual patches display heavy text-token gaps, rendering trend forecasts highly unstable:
\begin{equation}
\texttt{PatchedSignal}(?ps) \wedge \texttt{overlapRatio}(?ps, ?or) \wedge (?or < 0.5) \longrightarrow \texttt{exhibitsAliasing}(?ps, \texttt{H3\_Overlap})
\end{equation}